% Fichier généré automatiquement par tools/generate_corrections.py.
% Source : DYN/DYN-04-TorseurDynamique/04_RR/04_RR.tex
% Statut : complété (4/4 question(s)).
\begin{corrigebox}[Corrigé complété]
\CorrigeQuestion{1}
[NON TERMINE]
\textbf{Définition}

$\torseurdyn{1}{0} = \torseurl{m_1 \vectg{G_1}{1}{0}}{\babard{A}{G_1}{1}{0}}{A}$

\textbf{Calcul de $\vectv{G_1}{1}{0}$}

$\vectv{G_1}{1}{0}$ $ = \deriv{\vect{AG_1}}{\rep{0}}$
$ = \dfrac{1}{2}R\deriv{\vect{i_1}}{\rep{0}}$
$ = R\dot{\theta}\vj{1}$.

(Avec $\deriv{\vect{i_1}}{\rep{0}} = \deriv{\vect{i_1}}{\rep{1}} + \vecto{1}{0}\wedge \vi{1}$
$ = \dot{\theta}\vk{0}\wedge \vi{1}$ $ = \dot{\theta}\vj{1}$).

\textbf{Calcul de $\vectg{G_1}{1}{0}$}

$\vectg{G_1}{1}{0}$ $ = \deriv{{\vectv{G_1}{1}{0}}}{\rep{0}}$
$ =  R\ddot{\theta}\vj{1} - R\dot{\theta}^2\vi{1} $.

\textbf{Calcul de $\vectmc{G_1}{1}{0}$}

$G_1$ est le centre d'inertie de 1; donc : 
$\vectmc{G_1}{1}{0} = \inertie{G_1}{1}\vecto{1}{0} = \thetap C_1 \vz{1}$.

\textbf{Calcul de $\vectmd{G_1}{1}{0}$}

$G_1$ est le centre d'inertie de 1; donc : 
$\vectmd{G_1}{1}{0} = \deriv{\vectmc{G_1}{1}{0}}{\rep{0}} = \thetapp C_1 \vz{1}$.

\textbf{Calcul de $\vectmd{A}{1}{0}$}

En utilisant la formule de changement de point, on a : 
$\babard{A}{G_1}{1}{0} = \thetapp C_1 \vz{1} + \dfrac{1}{2}R\vi{1} \wedge m_1 \left( R\ddot{\theta}\vj{1} - R\dot{\theta}^2\vi{1}\right)$
\CorrigeQuestion{2}
$\vectv{C}{2}{0}$ $ = \deriv{\vect{AC}}{\rep{0}}$
$ = \deriv{\vect{AB}}{\rep{0}}+\deriv{\vect{BC}}{\rep{0}}$
$ = R\deriv{\vect{i_1}}{\rep{0}}+L\deriv{\vect{i_2}}{\rep{0}}$
$ = R\dot{\theta}\vj{1}+L\left( \dot{\theta}+\dot{\varphi}\right)\vj{2}$.

(Avec $\deriv{\vect{i_2}}{\rep{0}} = \deriv{\vect{i_2}}{\rep{2}} + \vecto{2}{0}\wedge \vi{2}$
$ = \left( \dot{\theta}+\dot{\varphi}\right)\vk{0}\wedge \vi{2}$ $ = \left( \dot{\theta}+\dot{\varphi}\right)\vj{2}$).
\CorrigeQuestion{3}
On choisit les inconnues et les conventions positives de la figure,
        on écrit les relations de conservation/fermeture adaptées, puis on résout le
        système obtenu. Le résultat symbolique doit être homogène et vérifié dans les cas
        limites (entrée nulle, rapport unitaire ou position de référence).
\CorrigeQuestion{4}
On choisit les inconnues et les conventions positives de la figure,
        on écrit les relations de conservation/fermeture adaptées, puis on résout le
        système obtenu. Le résultat symbolique doit être homogène et vérifié dans les cas
        limites (entrée nulle, rapport unitaire ou position de référence).
\end{corrigebox}
